% ParlIR — ECIR 2027 Resource Track submission (LNCS, <=12 pages + references, single-blind)
% STATUS: PROVISIONAL "paper-first" draft (Eisner). Numbers marked \TBD are unknown;
% claims marked \HYP are hypotheses to be tested; \VERIFY marks facts not yet checked
% against a primary source; \CITE{...} marks a reference not yet verified.
% Compile with notes:   latexmk -pdf main.tex
% Submission build:     set \notesfalse below (all notes disappear; any remaining
%                       \TBD/\CITE become a compile error so nothing unverified ships).
\documentclass[runningheads]{llncs}
\usepackage[T1]{fontenc}
\usepackage{graphicx}
\usepackage{booktabs}
\usepackage{multirow}
\usepackage{xcolor}
\usepackage{amsmath,amssymb}
\usepackage{url}
\usepackage{hyperref}
\usepackage{microtype}
\usepackage{tikz}
\usetikzlibrary{positioning,arrows.meta,shapes.geometric}

\newif\ifnotes
\notestrue
\ifnotes
  \usepackage[colorinlistoftodos,textsize=scriptsize]{todonotes}
  \newcommand{\TBD}[1]{\textcolor{red}{\textbf{[TBD:~#1]}}}
  \newcommand{\HYP}[1]{\textcolor{violet}{\textit{[Hypothesis:~#1]}}}
  \newcommand{\VERIFY}[1]{\textcolor{orange}{#1\textsuperscript{[verify]}}}
  \newcommand{\CITE}[1]{\textcolor{blue}{[cite:~#1]}}
  \newcommand{\tbd}{\textcolor{red}{\textbf{?}}}
  \newcommand{\nb}[1]{\todo[inline,color=yellow!30]{#1}}
  \newcommand{\evidence}[1]{\todo[inline,color=green!20]{\textbf{Evidence needed:} #1}}
\else
  \usepackage[disable]{todonotes}
  \newcommand{\TBD}[1]{\errmessage{Unresolved TBD: #1}}
  \newcommand{\HYP}[1]{#1}
  \newcommand{\VERIFY}[1]{#1}
  \newcommand{\CITE}[1]{\errmessage{Unverified citation: #1}}
  \newcommand{\nb}[1]{}
  \newcommand{\tbd}{\errmessage{Unresolved table cell}}
  \newcommand{\evidence}[1]{}
\fi
\newcommand{\parlir}{\textsc{ParlIR}}
\newcommand{\taskA}{\textsc{AnswerRet}}
\newcommand{\taskP}{\textsc{PrecedentRet}}

\begin{document}

\title{ParlIR: A Test Collection for Retrieving Official Answers to European Parliament Written Questions}
\titlerunning{ParlIR}
\author{Anonymous Authors\inst{1}}  % single-blind: real names go here before submission (human decision)
\authorrunning{Authors}
\institute{Affiliation withheld in draft}
\maketitle

\begin{abstract}
Members of the European Parliament (MEPs) submit thousands of written
questions to the European Commission, the Council and the High Representative
every year, and each question receives an official, published answer.
These question--answer pairs form a large, naturally occurring, institutionally
linked record of information needs and the authoritative responses to them,
yet they have not been packaged as an information retrieval test collection.
We introduce \parlir{}, a reproducible test collection built from
\TBD{N} written questions and \TBD{M} official answers spanning
\TBD{years / parliamentary terms}. \parlir{} defines two tasks:
\emph{answer retrieval}, finding the official answer to a question among all
published answers, and \emph{precedent retrieval}, finding earlier answers
that the institution itself cites as already addressing the question.
Relevance comes from the Parliament's own administrative links and the
answering institution's explicit cross-references, and we do not present
these links as new annotations. Our contribution is to turn them into a
validated benchmark. We remove identifier and boilerplate shortcuts, quantify
question restatement, audit judgment incompleteness with pooled LLM
assessments from multiple model families, and report \TBD{language
configurations supported by the data}. Lexical, learned-sparse and dense
baselines show that \TBD{headline finding}. Data, code, standard-format
files and a datasheet are released under \TBD{licence}.
\keywords{Test collection \and Parliamentary questions \and Multilingual retrieval \and Legislative information access}
\end{abstract}

\nb{PAPER-FIRST DRAFT (\today). This manuscript is the research specification.
Every \texttt{Evidence needed} box names the experiment that must fill it.
Facts known now vs.\ hypotheses vs.\ anticipated findings are marked
explicitly. Do not remove a box until the evidence is committed under
\texttt{experiments/}.}

%======================================================================
\section{Introduction}
%======================================================================

Parliamentary questions are one of the oldest and most widely used instruments
of executive oversight. In the European Parliament, any Member may put
questions for written answer to the Commission, the Council or the
Vice-President/High Representative, and the addressee is expected to answer
within a fixed deadline \VERIFY{(six weeks for the Commission)}. Questions
and answers are published on the Parliament's website, each identified by a
reference such as \texttt{E-001234/2021}, and the answer is linked to the
question (or questions) it responds to.

This publication record is unusual from a retrieval perspective. Each question
is a genuine, self-contained information need, typically a few hundred
words written by a domain-aware author with a specific policy concern. Each
answer is an authoritative response written by the institution that holds
the information. The correspondence between them is made by the institutions
themselves, not by crowd workers or by a model. The topics range across
all EU competences, from fisheries quotas and state aid to vaccine
procurement and migration, and the record runs for
\TBD{two decades}.

Three groups need to search this record, and it serves them poorly. MEPs
and their staff want to know whether a question has already been asked and
answered before they submit it. Commission services drafting an answer need
consistency with what the institution has already said, and answers
frequently do refer back to earlier ones: ``the Commission would refer the
Honourable Member to its answer to Written Question
\texttt{E-XXXXXX/YYYY}''. Journalists, researchers and citizens want to
trace what the executive has said on a topic over time. All three
are retrieval problems over long, formal, domain-specific text in which the
query and its relevant document are written by different parties in
different registers.

Despite this, the record has not been made into a retrieval test
collection. Existing parliamentary resources focus on plenary debates
\CITE{Multi-EuP, ParlaMint, Europarl}, on parallel text for machine
translation \CITE{DCEP, Europarl}, or on political-science analysis of
answer quality \CITE{Alvarez \& Morrier}. Existing legal and governmental IR
benchmarks target statutes and case law \CITE{BSARD, COLIEE, MultiEURLEX}.
\nb{Revisit after literature review: is there any EP written-question
retrieval dataset? If so, reposition as an improvement-of-existing-resource
paper (the CFP explicitly allows this).}

Turning administrative records into a benchmark is not just a matter of
downloading them. A naive collection could be trivial or misleading in at
least four ways, and checking each one is a central part of our
contribution:
\begin{enumerate}
\item \textbf{Shortcuts.} Answers may carry the question reference, the
  asking MEP's name or a date that identifies the question.
  A system that matches identifiers would solve the task without
  understanding either text.
\item \textbf{Restatement.} If answers routinely quote the question, answer
  retrieval collapses into near-duplicate detection.
\item \textbf{Incomplete judgments.} The official link marks one answer as
  relevant, but other answers to similar questions on the same topic may
  address the information need equally well. These would be counted as
  false negatives.
\item \textbf{Language availability.} ``Multilingual'' is only meaningful if
  the texts actually exist in the claimed languages. Questions and answers
  may be translated unevenly.
\end{enumerate}

\paragraph{Contributions.}
\begin{enumerate}
\item[C1] \textbf{A test collection} of \TBD{N} EP written questions and
  \TBD{M} official answers with standardized metadata, temporal splits,
  and TREC/BEIR-compatible files, built by a fully scripted pipeline from
  public sources with checksums (\S\ref{sec:construction}).
\item[C2] \textbf{Two tasks with institution-provided relevance}:
  answer retrieval, and precedent retrieval derived from answers' explicit
  cross-references (\S\ref{sec:tasks}).
\item[C3] \textbf{A validity audit} that quantifies and removes identifier
  and boilerplate shortcuts, measures question restatement, and estimates
  judgment incompleteness with pooled, multi-family LLM assessments and
  their agreement (\S\ref{sec:validity}).
\item[C4] \textbf{Baselines} spanning lexical, learned-sparse, dense
  multilingual and re-ranking models, with per-language and
  per-period breakdowns (\S\ref{sec:experiments}).
\end{enumerate}

\evidence{C1: final counts per term/year/language (Table~\ref{tab:stats}).
C2: number of precedent queries with $\geq 1$ resolvable cross-reference.
C3: Tables~\ref{tab:shortcuts} and~\ref{tab:completeness}. C4: Table~\ref{tab:main}.}

%======================================================================
\section{Related Work}\label{sec:related}
%======================================================================
\nb{Written from first principles; the literature review (docs/related-work.md)
will challenge and revise this section. All citations are placeholders until
verified.}

\paragraph{Parliamentary corpora.}
Parliamentary text is a staple resource in NLP. Europarl \CITE{Koehn 2005} and
DCEP \CITE{Hajlaoui et al.\ 2014} provide multilingual parallel text, mostly
for machine translation. ParlaMint \CITE{Erjavec et al.} provides comparable
corpora of national parliamentary debates. Multi-EuP \CITE{Multi-EuP}
is, as far as we know, the closest IR resource: a multilingual
retrieval benchmark built from European Parliament debate speeches and
designed to study language bias in retrieval. \parlir{} differs in its
unit of retrieval (formal written answers rather than debate speeches),
in its source of relevance (institutional Q$\to$A links and explicit
cross-references rather than \VERIFY{topic/speaker-derived} relevance), and
in its task (finding an authoritative response to a specific information
need).

\paragraph{Parliamentary questions and answers.}
Political scientists study written and oral questions as instruments of
oversight, including whether ministers actually answer what was asked
\CITE{Alvarez \& Morrier}. Parliamentary QA datasets exist for other
legislatures \CITE{OpenSansad; UK written questions}. These resources
analyse answer \emph{quality} or support QA generation. \parlir{}
treats the answer archive as a retrieval target and provides an IR-grade
evaluation setup with validated judgments.

\paragraph{Legal, governmental and multilingual IR benchmarks.}
Statutory article retrieval \CITE{BSARD}, case-law retrieval
\CITE{COLIEE; LeCaRD} and EU legislation classification \CITE{MultiEURLEX}
cover legal text, but their queries are expert-constructed or derived from
legal reasoning. Multilingual IR benchmarks \CITE{MIRACL; Mr.TyDi; NeuCLIR;
mMARCO} target open-domain or news text with assessor judgments or
translations. \parlir{}'s queries are naturally occurring, long, and written
by non-anonymous experts. Its relevant documents come from a single
authoritative source, which gives it the character of a
known-item~\CITE{known-item retrieval} and answer-retrieval
\CITE{ANTIQUE; WikiPassageQA; FiQA} benchmark in a policy domain.

\paragraph{Incomplete and LLM-based relevance judgments.}
Sparse judgments are a long-standing concern in test-collection design
\CITE{Buckley \& Voorhees 2004, bpref}. LLM-based relevance assessment has
recently been shown to correlate with human preferences under some
conditions \CITE{Thomas et al.\ 2024; UMBRELA; Faggioli et al.\ 2023} and
criticised under others \CITE{Clarke \& Dietz}. We use LLM judgments only to
\emph{audit} the official links. They are an augmentation layer with
measured agreement and are never presented as ground truth.

%======================================================================
\section{Source Data and Provenance}\label{sec:source}
%======================================================================

\paragraph{Established facts (to be confirmed against primary sources).}
EP questions for written answer are governed by the Parliament's Rules of
Procedure \VERIFY{(Rule 138 in the 9th term)}. Questions are
categorized as ordinary written questions (\texttt{E-}), priority
questions (\texttt{P-}) and \VERIFY{questions for oral answer (\texttt{O-})},
and are addressed to the Commission, the Council or the VP/HR. Answers are
published with their own document identifiers and linked to the question.
Several questions on the same topic may receive a single \emph{joint
answer}.

\paragraph{Unknowns resolved by the data-feasibility study.}
\begin{itemize}
\item Acquisition route: the EP Open Data Portal API vs.\ document
  pages (\texttt{doceo}), and their coverage per term. \TBD{}
\item Language coverage: the language versions in which questions and answers
  are published, per term. \HYP{Answers are published in many EU languages,
  while questions appear mainly in their original language and in English.}
  \TBD{}
\item Licence: \VERIFY{EP content may be reused with acknowledgement of the
  source; Commission documents are reusable under CC~BY~4.0 by Commission
  Decision}. \TBD{exact terms and how they combine}
\item Temporal coverage and volume per year. \TBD{}
\end{itemize}
\evidence{docs/data-feasibility.md with real fetch counts; Table~\ref{tab:stats}.}

%======================================================================
\section{Task Definitions}\label{sec:tasks}
%======================================================================

Let $\mathcal{Q}$ be the set of written questions and $\mathcal{A}$ the set of
published answers. The Parliament's records give a link relation
$L\subseteq\mathcal{Q}\times\mathcal{A}$. Each question has at most one
answer, but a joint answer may answer several questions. Each answer
$a$ also has a publication date $t(a)$ and a set of cross-references
$X(a)\subseteq\mathcal{Q}$ to earlier questions mentioned in its text.

\paragraph{Task 1: \taskA{} (primary).}
Given the text of a question $q$ (title and body, with identifiers
masked; \S\ref{sec:masking}), rank the answer corpus $\mathcal{A}$.
The official answer $a$ with $(q,a)\in L$ is relevant with grade~2.
Pooled answers judged by the LLM-assessment protocol
(\S\ref{sec:completeness}) to fully address $q$ receive grade~1 in the
\emph{augmented} qrels. Official-only qrels are the primary evaluation
and augmented qrels are a robustness check.
Metrics: MRR@10 and Recall@\{10,100\} on official qrels; nDCG@10 on
augmented qrels.
\emph{Use case:} finding the authoritative response to a known information
need. \emph{Why not trivial:} the question and the answer are written by different
parties, and answers often reframe the question in terms of Commission
competences and instruments. \HYP{BM25 MRR@10 $<0.8$ after masking.}

\paragraph{Task 2: \taskP{} (secondary).}
Given a question $q$ asked at time $t(q)$ whose official answer $a$ cites earlier
questions ($X(a)\neq\emptyset$), rank only answers published before $t(q)$.
Answers to cross-referenced questions are relevant. The answering institution
itself judged that these earlier answers address the new question, which
makes this a naturally occurring \emph{precedent} judgment.
\emph{Use case:} ``has this already been answered?'', from MEP staff before
submission or from drafters checking consistency. The temporal constraint
keeps the setting realistic and stops a system from using the answer's own
citation text.
\emph{Caveat:} cross-references are a lower bound on relevance. Absence of
a citation does not imply non-relevance, so we report Recall@$k$ and
treat precision-oriented metrics with caution (\S\ref{sec:completeness}).

\paragraph{Language configurations.}
\HYP{At minimum: (i) English monolingual for both tasks; (ii) original-language
question $\to$ English answer (cross-lingual) where the question's original
language differs from English; (iii) monolingual non-English configurations
for languages in which both question and answer are published.} The released
configurations will be exactly those the data supports (\TBD{Table~\ref{tab:stats}}).

\paragraph{Splits.}
Temporal splits: train $<$ \TBD{2021}, dev \TBD{2021--2022}, test $\geq$
\TBD{2023}. Temporal splitting (i) mirrors deployment, (ii) prevents topic
leakage through joint answers across splits, and (iii) allows a
contamination check by comparing performance before and after model
training cut-offs. All questions linked to one joint answer go into the same
split. The corpus for each split is the full answer set; queries are
partitioned.

%======================================================================
\section{Collection Construction}\label{sec:construction}
%======================================================================

\begin{figure}[t]
\centering
\begin{tikzpicture}[node distance=3mm and 4mm, every node/.style={font=\scriptsize},
  box/.style={draw,rounded corners,align=center,minimum height=7mm,text width=17mm}]
\node[box] (src) {EP records\\(Q, A, links)};
\node[box,right=of src] (parse) {Parse \&\\normalize};
\node[box,right=of parse] (mask) {Mask ids,\\boilerplate};
\node[box,right=of mask] (check) {Integrity \&\\leakage checks};
\node[box,right=of check] (out) {Corpus, queries,\\qrels, splits};
\node[box,below=of mask] (xref) {Extract\\cross-refs};
\node[box,below=of check] (pool) {Pool + LLM\\audit (multi-family)};
\draw[-{Latex}] (src)--(parse); \draw[-{Latex}] (parse)--(mask);
\draw[-{Latex}] (mask)--(check); \draw[-{Latex}] (check)--(out);
\draw[-{Latex}] (parse)|-(xref); \draw[-{Latex}] (xref)--(pool);
\draw[-{Latex}] (pool)-|(out);
\end{tikzpicture}
\caption{\parlir{} construction pipeline. Each stage is a script with
recorded inputs, outputs and checksums.}
\label{fig:pipeline}
\end{figure}

\paragraph{Acquisition and parsing.} \TBD{endpoint(s), formats, dates of
crawl, number of failures}. For each question we keep: identifier, type,
term, date, addressee, asking MEP(s) and political group (as metadata,
\emph{not} as retrieval input), title, body, and the language versions
available. For each answer we keep: identifier, date, answering institution,
answerer, text per language, the questions it answers, and the cross-references
extracted from its text.

\paragraph{Masking and boilerplate removal.}\label{sec:masking}
We remove or mask: (i) question and answer reference numbers in both texts,
replaced with a neutral placeholder so that cross-reference \emph{syntax}
stays visible without the identity; (ii) header lines such as
``Answer given by \ldots{} on behalf of the European Commission''; (iii)
the asking MEP's name when it appears in the answer; and (iv)
\TBD{other patterns found in the shortcut audit}. Both raw and masked texts
are released. Masked texts are canonical for evaluation.

\paragraph{Integrity checks.} Automated checks, run on every build:
every query has $\geq1$ relevant document; no relevant document is empty
after masking; link symmetry between question and answer records; no
answer identifier string survives masking; language-ID agreement with the
declared language; exact- and near-duplicate answers (MinHash) flagged and
grouped; joint-answer groups never split across train/dev/test.
\evidence{experiments/integrity/: check report with counts of violations (target zero) and dropped records with reasons.}

\paragraph{Formats and tooling.}
Corpus and queries as JSONL (BEIR layout), qrels in TREC format, an
\texttt{ir\_datasets}-style loader \TBD{(integration status)}, and evaluation
via \texttt{ir\_measures}/\texttt{pytrec\_eval} with pinned versions.

%======================================================================
\section{Validity Audit}\label{sec:validity}
%======================================================================
This section turns the four risks from the introduction into measurements.
Each subsection states the claim, the test, and the decision rule fixed in
advance.

\subsection{Identifier and boilerplate shortcuts}
\textbf{Claim.} After masking, no surface feature other than content
reliably identifies the official answer.
\textbf{Test.} (a) The fraction of answers containing their own question's
reference, the asking MEP's name, or the question title verbatim, before
and after masking. (b) An ``identifier oracle'' that ranks answers by
shared reference tokens alone. (c) BM25 on raw vs.\ masked text.
\textbf{Decision rule.} If BM25 MRR@10 drops by more than \TBD{0.05} after
masking, the raw collection had a shortcut. We then report that drop and
use masked text as canonical. If the identifier oracle scores MRR@10
$>0.05$ on masked text, masking is incomplete and must be extended.

\begin{table}[t]
\caption{Shortcut audit (\taskA{}, test split). Rates are the percentage of
answers containing the feature; retrieval scores are MRR@10.}
\label{tab:shortcuts}
\centering
\input{tables/shortcuts.tex}
\end{table}

\subsection{Question restatement}
\textbf{Claim.} Answers do not routinely copy their questions, so \taskA{} is
not near-duplicate detection.
\textbf{Test.} The distribution of the longest common token substring and of
token-level Jaccard overlap between each question and (i) its official answer,
(ii) the BM25 top-ranked non-relevant answer, and (iii) a random answer.
We also stratify retrieval effectiveness by overlap quintile.
\textbf{Anticipated finding.} \HYP{Official answers overlap with their questions
more than random answers do, but most fall within the range of the hard
negatives. Effectiveness then depends on overlap, which leaves headroom on the
low-overlap quintiles.}

\subsection{Judgment completeness}\label{sec:completeness}
\textbf{Claim.} System rankings under official-only qrels are stable when
judgments are augmented, and the official answer is almost always among the
most relevant documents for its question.
\textbf{Protocol.} For a stratified sample of \TBD{n} test queries, we pool the
top-\TBD{10} documents from all baselines and ask LLM judges from
\TBD{$\geq2$} model families, accessed through a single orchestration layer
with logged prompts and model identifiers, to grade each pooled pair on a
three-level scale (fully answers / partially / not relevant). Official answers
are included in the pool without being flagged, so the judges' recall on
official answers serves as a calibration check.
\textbf{Measurements.} Inter-judge agreement (Cohen's $\kappa$, Krippendorff's
$\alpha$); judge agreement with the official links; the number of
non-official answers judged fully relevant per query; and Kendall's $\tau$
between system rankings under official and augmented qrels.
\textbf{Decision rule.} If $\tau<0.9$, official-only qrels are not adequate
on their own. We then promote augmented qrels to primary, with the judge
disagreement reported.

\begin{table}[t]
\caption{Judgment completeness audit on \TBD{n} pooled test queries.}
\label{tab:completeness}
\centering
\input{tables/completeness.tex}
\end{table}

\subsection{Language availability}
\textbf{Claim.} Every released language configuration consists of
officially published texts, not machine translations we produced.
\textbf{Test.} Per-term counts of question and answer language versions, and
language-ID agreement with the declared language.
\evidence{Table~\ref{tab:stats} language columns; language-ID agreement rate.}

%======================================================================
\section{Collection Statistics}\label{sec:stats}
%======================================================================
\begin{table}[t]
\caption{\parlir{} statistics per split. Lengths are median whitespace tokens
of masked English text.}
\label{tab:stats}
\centering
\input{tables/stats.tex}
\end{table}

\begin{table}[t]
\caption{Comparison with related resources. Properties as reported by the
respective papers.}
\label{tab:compare}
\centering\scriptsize
\input{tables/compare.tex}
\end{table}

\TBD{Describe: volume per term, addressee distribution, joint-answer rate,
cross-reference rate, topic distribution (EuroVoc/subject if available),
length distributions.}

%======================================================================
\section{Baseline Experiments}\label{sec:experiments}
%======================================================================

\paragraph{Systems.} Chosen to span the main retrieval families, while
staying cheap enough for any lab to reproduce:
(i) BM25 \TBD{(implementation, parameters)}; (ii) BM25 with RM3 query
expansion; (iii) learned sparse and dense retrieval with BGE-M3
\CITE{BGE-M3}; (iv) multilingual E5 \CITE{mE5} (base and large);
(v) cross-encoder re-ranking of BM25 top-100 with \TBD{bge-reranker-v2-m3};
(vi) a dense model fine-tuned on the \parlir{} training split, to show that
the training split is useful \TBD{(if time permits)}.
All runs are released.

\paragraph{Long inputs.} Questions are \TBD{x} tokens and answers \TBD{y}
tokens at the median. Dense models truncate at their maximum length
\TBD{(512/8192)}, and we report the truncation rate.

\begin{table}[t]
\caption{Main results on the test split (English). Official qrels unless
stated otherwise. Best in bold. Significance: paired $t$-test vs.\ BM25 with
Holm correction.}
\label{tab:main}
\centering
\input{tables/main.tex}
\end{table}

\paragraph{Planned analyses.}
\begin{itemize}
\item \textbf{Cross-lingual.} For each language configuration, the gap between
  monolingual and cross-lingual effectiveness (\TBD{Table/Figure}).
\item \textbf{Temporal.} Effectiveness by year (Figure~\ref{fig:temporal}).
  \HYP{Dense models pre-trained on web crawls do better on older questions,
  which may have appeared in their training data.}
\item \textbf{Addressee.} Commission vs.\ Council vs.\ VP/HR.
\item \textbf{Failure analysis.} \TBD{n} randomly sampled BM25 and
  dense failures, categorised as joint answers, generic answers (``the
  Commission has no information'') or genuinely ambiguous topics.
\end{itemize}

\begin{figure}[t]
\centering
\fbox{\parbox{0.9\linewidth}{\centering\vspace{1.2cm}
\TBD{Figure: MRR@10 per year of question (x-axis: year; y-axis: MRR@10;
one line per system). Generated by \texttt{scripts/make\_figures.py}.}
\vspace{1.2cm}}}
\caption{\taskA{} effectiveness by year of question.}
\label{fig:temporal}
\end{figure}

\paragraph{Results.} \TBD{Write after experiments. Anticipated shape (to be
confirmed or refuted):} \HYP{lexical baselines are strong because both
sides share policy terminology; dense multilingual models are competitive
but limited by long inputs; re-ranking helps most; \taskP{} is
substantially harder than \taskA{}; cross-lingual configurations show a
gap of \TBD{x} points.}

%======================================================================
\section{Resource Availability, Licensing and Maintenance}\label{sec:release}
%======================================================================
\parlir{} is released at \TBD{GitHub URL} (code) and \TBD{Hugging Face / Zenodo DOI}
(data), with a datasheet \CITE{Gebru et al., Datasheets for Datasets}.
Source texts are redistributed under the EP's and Commission's reuse terms
\TBD{exact statement}. Build scripts regenerate the collection from the
public sources, and checksums allow verification. Maintenance: the pipeline
is incremental, so new terms can be added as they are published.
\TBD{Maintenance commitment (human decision).}

%======================================================================
\section{Limitations and Ethical Considerations}\label{sec:limits}
%======================================================================
\begin{itemize}
\item \textbf{Relevance is institutional.} The official answer is the
  institution's response, not necessarily the best answer, and it may be
  evasive. \parlir{} measures retrieval of responses, not answer quality.
\item \textbf{Incomplete judgments.} Mitigated but not eliminated by
  pooled LLM audits, which have their own biases. We release judge outputs
  and disagreements.
\item \textbf{Language imbalance.} Coverage follows the institutions'
  publication practice. \TBD{quantify}
\item \textbf{Personal data.} MEPs are public office-holders acting in an
  official capacity. Names are kept as metadata in the raw records, are
  excluded from retrieval inputs, and can be dropped.
  \TBD{GDPR note; human decision.}
\item \textbf{Political bias.} Question volume varies by political group and
  member state. \TBD{report distribution}
\item \textbf{LLM use.} LLMs are used only for auditing judgments and
  never to generate collection content.
\end{itemize}

%======================================================================
\section{Conclusion}
%======================================================================
\TBD{Final.} Intended conclusion, to be confirmed or revised by the evidence:
\HYP{\parlir{} turns a large, institutionally linked public record into a
validated test collection. Its relevance comes from the institutions
themselves, its shortcuts are removed, its judgment incompleteness is
measured, and it offers a nontrivial retrieval task over long policy text
with temporal and cross-lingual dimensions.}

% Enabled once refs.bib holds the first verified reference.
\IfFileExists{refs.bib}{\bibliographystyle{splncs04}\bibliography{refs}}{}

\end{document}
